\documentclass[11pt]{article}
\usepackage[margin=1in]{geometry}
\usepackage[T1]{fontenc}
\usepackage[utf8]{inputenc}
\usepackage{amsmath,amssymb,amsthm}
\usepackage{enumitem}

\title{ICPC World Finals 2020\\F. Ley Lines}
\author{}
\date{}

\begin{document}
\maketitle

\section*{Problem Summary}

We are given $n$ planar points and a pencil thickness $t$.
A ``line'' of thickness $t$ is a strip between two parallel lines whose distance is $t$.
We must find the maximum number of points that can lie inside one such strip.

\section*{Key Observations}

\begin{itemize}[leftmargin=*]
    \item Any optimal strip can be translated orthogonally until one of its two boundary lines passes through
    one of the chosen points. So it is enough to consider strips whose lower boundary passes through some
    input point $P$.
    \item Fix such a point $P$ and let the boundary line through $P$ have direction angle $\alpha$.
    For another point $Q$, write:
    \begin{itemize}[leftmargin=*]
        \item $d = |PQ|$,
        \item $\gamma =$ the polar angle of vector $\overrightarrow{PQ}$.
    \end{itemize}
    Then $Q$ lies inside the strip exactly when its signed perpendicular distance from the boundary line is
    between $0$ and $t$, i.e.
    \[
        0 \le d \sin(\gamma - \alpha) \le t.
    \]
    \item Therefore, for fixed $P$, every other point contributes one or two angle intervals of valid
    $\alpha$ values.
    \item If $d \le t$, then any line direction that keeps $Q$ on the chosen side of the boundary works, so the
    valid interval has length $\pi$:
    \[
        \alpha \in [\gamma-\pi,\ \gamma].
    \]
    \item If $d > t$, let $\delta = \arcsin(t/d)$.
    Then $Q$ is inside the strip exactly for
    \[
        \alpha \in [\gamma-\delta,\ \gamma]
    \]
    or
    \[
        \alpha \in [\gamma-\pi,\ \gamma-\pi+\delta].
    \]
    \item So for each anchor point $P$, the problem becomes: on the angle circle $[0,2\pi)$, find the maximum
    number of these intervals that overlap at one angle.
\end{itemize}

\section*{Algorithm}

\begin{enumerate}[leftmargin=*]
    \item Iterate over every point $P_i$ and force one strip boundary to pass through it.
    \item For every other point $P_j$:
    \begin{itemize}[leftmargin=*]
        \item compute $d=|P_iP_j|$ and $\gamma=\operatorname{atan2}(y_j-y_i, x_j-x_i)$;
        \item convert the valid-angle set described above into one or two intervals on the angle circle;
        \item split wrap-around intervals at $0$ if necessary.
    \end{itemize}
    \item Run a sweep over the interval endpoints:
    \begin{itemize}[leftmargin=*]
        \item start event = $+1$,
        \item end event = $-1$,
        \item when two events have the same angle, process starts before ends because interval endpoints are
        inclusive.
    \end{itemize}
    \item The sweep count plus the anchor point itself gives the best strip with boundary through $P_i$.
    \item Take the maximum over all anchors.
\end{enumerate}

\section*{Correctness Proof}

We prove that the algorithm returns the correct answer.

\paragraph{Lemma 1.}
There exists an optimal strip whose one boundary line passes through an input point.

\paragraph{Proof.}
Take any optimal strip and project its chosen points onto the strip normal direction.
All projected values lie inside some interval of length $t$.
Shift that interval until its left endpoint equals the minimum projected value among those chosen points.
This translation does not remove any chosen point, and now the left boundary passes through at least one of
them.
So some optimal strip has a boundary through an input point. \qed

\paragraph{Lemma 2.}
Fix an anchor point $P$ on the lower boundary of the strip, and let the boundary line have direction angle
$\alpha$.
For another point $Q$ at polar angle $\gamma$ and distance $d$ from $P$, the set of valid $\alpha$ is exactly:
\[
    [\gamma-\pi,\gamma] \quad \text{if } d \le t,
\]
and otherwise, with $\delta=\arcsin(t/d)$,
\[
    [\gamma-\delta,\gamma] \cup [\gamma-\pi,\gamma-\pi+\delta].
\]

\paragraph{Proof.}
Point $Q$ is inside the strip exactly when its signed perpendicular distance from the boundary line is
between $0$ and $t$:
\[
    0 \le d\sin(\gamma-\alpha) \le t.
\]

If $d \le t$, the upper bound is automatic, so we only need
\[
    \sin(\gamma-\alpha) \ge 0,
\]
which is equivalent to $\gamma-\alpha \in [0,\pi]$, or
\[
    \alpha \in [\gamma-\pi,\gamma].
\]

If $d > t$, the inequality becomes
\[
    0 \le \sin(\gamma-\alpha) \le t/d.
\]
Writing $\delta=\arcsin(t/d)$, this holds exactly when
\[
    \gamma-\alpha \in [0,\delta] \cup [\pi-\delta,\pi],
\]
which is equivalent to the two stated intervals for $\alpha$. \qed

\paragraph{Lemma 3.}
For a fixed anchor point $P$, the sweep over interval endpoints computes the maximum number of points that
can be covered by a strip whose lower boundary passes through $P$.

\paragraph{Proof.}
By Lemma 2, each other point $Q$ contributes exactly the set of boundary angles $\alpha$ for which $Q$ lies
inside the strip.
Hence, for any specific angle $\alpha$, the number of points covered equals one (for the anchor $P$) plus
the number of intervals containing $\alpha$.
The interval-endpoint sweep computes the maximum overlap of these intervals on the angle circle, so it
computes exactly the best strip anchored at $P$. \qed

\paragraph{Theorem.}
The algorithm outputs the correct answer for every valid input.

\paragraph{Proof.}
By Lemma 1, some optimal strip has a boundary through an input point $P$.
For that fixed anchor, Lemma 3 shows that the algorithm computes the best possible strip.
Since the algorithm tries every point as anchor, it will examine the anchor from some optimal strip and
therefore reach the global optimum.
Thus the reported maximum is correct. \qed

\section*{Complexity Analysis}

For each of the $n$ anchor points we generate $O(n)$ intervals and sort $O(n)$ events.
So the total running time is
\[
    O(n^2 \log n).
\]

The event list for one anchor uses $O(n)$ memory.

\section*{Implementation Notes}

\begin{itemize}[leftmargin=*]
    \item Wrapped intervals on the circle are split into two ordinary intervals in $[0,2\pi)$.
    \item Endpoints are treated as inclusive, so when two events share the same angle we process start events
    before end events.
    \item \texttt{long double} is sufficient; the statement guarantees the answer is stable under a
    $10^{-2}$ perturbation of $t$.
\end{itemize}

\end{document}
